% ============================================================================
% 00_intro.tex -- Introduction (~1/2 page)
% All AI-drafted prose is wrapped in \ai{...}; revise personally before submission.
% ============================================================================
\section{Introduction}
\label{sec:intro}

\ai{The one-dimensional Kitaev chain supports boundary Majorana modes in its
topological phase~\cite{kitaev2001}.  After a Jordan--Wigner transformation,
however, neither the Majorana wave functions nor an exact eigenvector are direct
hardware observables.  The relevant experimental question is therefore more
operational: can a parameterized quantum circuit prepare the finite-chain ground
state and can measurements of that circuit distinguish the topological and
trivial regimes?}

\ai{This report covers Part~4 of the project, \emph{Measuring Topology on
Circuits}.  The qubit Hamiltonian and parity convention are taken as inputs from
J. Singh's Part~3 report rather than re-derived here.  I construct the ideal
circuit workflow: ground-state preparation with a variational quantum
eigensolver (VQE), comparison of local and non-local observables, basis-rotated
shot readout of the edge string, a chemical-potential sweep, parity-constrained
optimization, and an ansatz-depth diagnostic.  The observable used below is a
finite-size boundary/string diagnostic; it is not, by itself, a replacement for
the bulk winding number.  Gate, relaxation, and readout-noise effects are outside
the present scope and are treated in D. Stoev's Part~5 report.}
